%% kapa: exam=1 round=35 year=2024 period=2 form=A subject=appraisal_law subject_ko=감정평가관계법규 seq=37 number=37
%% source: Q-Net 감정평가사 기출문제 (cst00309) artlSeq=5213229 file=raw/1008/5213229/2교시 문제지 원본.pdf
%% conversion: kordoc 4.24.51 -> tools/gen1.py (auto); review=proofread-sonnet; flags=none
\begin{question}{37}
도시 및 주거환경정비법령상 정비사업에 관한 설명으로 옳은 것은?
\begin{choices}
  \item 재개발사업이란 정비기반시설은 양호하나 노후ㆍ불량건축물에 해당하는 공동주택이 밀집한 지역에서 주거환경을 개선하기 위한 사업을 말한다.
  \item 재건축사업의 경우 정비구역의 지정권자는 하나의 정비구역을 둘 이상의 정비구역으로 분할하여 지정할 수 없다.
  \item 재건축사업은 관리처분계획에 따라 건축물을 공급하거나 환지로 공급하는 방법으로 한다.
  \item 재개발사업의 시행자가 작성하는 사업시행계획서에는 임시거주시설을 포함한 주민이 주대책이 포함되어야 한다.
  \item 토지등소유자가 20인 미만인 경우에는 토지등소유자가 직접 재개발사업을 시행할 수 없다.
\end{choices}
\end{question}
\answer{4}
